\Page{8}{\textenglish{02 / FORECAST}}{מה נלבש בעונה הבאה}{48}{55}

\FeatureIntro{PaperSteel}{קניינית מסתכלת על לוח צבע לעונה שעוד רחוקה. היא יכולה להציע יותר פריטים בגוון שמוכר לקהל, או לתת מקום לצבע שמופיע שוב ושוב בתצלומים מהרחוב. הבחירה תופסת מקום אמיתי. צריך להזמין חומר, להקצות מקום על המדף ולהשאיר תקציב לצבעים אחרים. חזאית אופנה עוזרת לקרוא את הסימנים לפני שההחלטות האלה מתקבעות.}

\vspace{6mm}

{\fontsize{16}{25}\selectfont

העבודה המסורתית דרשה ללכת אל המקומות שבהם טעם מתגבש. חזאיות סיירו בירידים, ראו תצוגות, עקבו אחרי חנויות ותרבות רחוב והשוו את מה שראו לעונות קודמות. ג׳וסלין קורנר מ־\textenglish{WGSN} סיפרה שכיועצת צעירה הייתה סופרת פריטים על מתלים ומרכזת מחירים ומבחר בגיליון. ברגע שהרשומות מצטברות, אפשר להתחיל להבדיל בין פריט מפתיע שהופיע בחנות אחת לבין החלטה שחוזרת אצל כמה מותגים. השאלה היא אם הדמיון הזה משמעותי גם לקהל שעבורו מתכננים את העונה. \SourceRef{4}

שימוש בחזאות משתנה לפי ההחלטה. מעצבת שבונה קולקציה מבקשת כיוון לצבעים, לבדים ולצלליות. קניינית בודקת אם להרחיב קבוצה מסוימת ומה יהיה מועד ההשקה. מנהלת מסחר עשויה להתעניין בכמות שכדאי להתחייב לה לעומת כמות שאפשר להשלים אם יש ביקוש. כל אחת משתמשת באותה תצפית בקנה מידה אחר. \SourceRef{4}\SourceRef{5}

קשה לנתק את המקצוע מהזמן שבו אפשר לפעול. אם לוקח חודשים להזמין ולייצר, מידע מעונה שכבר בשיאה עלול להגיע באיחור. אם יש חומר זמין ואפשר להרחיב גוון של דגם מוכר, גם אות מצומצם יכול להצדיק ניסיון. המיומנות היא לזהות מה שווה החלטה עכשיו, ואחרי מה כדאי להמשיך לעקוב לפני שמתחייבים.

\par}
